\UseRawInputEncoding
\documentclass[a4paper, 14pt]{extarticle}

% Поля
%--------------------------------------
\usepackage{geometry}
\geometry{a4paper,tmargin=2cm,bmargin=2cm,lmargin=3cm,rmargin=1cm}
%--------------------------------------


%Russian-specific packages
%--------------------------------------
\usepackage[T2A]{fontenc}
\usepackage[utf8]{inputenc} 
\usepackage[english, main=russian]{babel}
%--------------------------------------

\usepackage{textcomp}

% Красная строка
%--------------------------------------
\usepackage{indentfirst}               
%--------------------------------------             


%Graphics
%--------------------------------------
\usepackage{graphicx}
\graphicspath{ {./images/} }
\usepackage{wrapfig}
%--------------------------------------

% Полуторный интервал
%--------------------------------------
\linespread{1.3}                    
%--------------------------------------

%Выравнивание и переносы
%--------------------------------------
% Избавляемся от переполнений
\sloppy
% Запрещаем разрыв страницы после первой строки абзаца
\clubpenalty=10000
% Запрещаем разрыв страницы после последней строки абзаца
\widowpenalty=10000
%--------------------------------------

%Списки
\usepackage{enumitem}

%Подписи
\usepackage{caption} 

%Гиперссылки
\usepackage{hyperref}

\hypersetup {
	unicode=true
}

%Рисунки
%--------------------------------------
\DeclareCaptionLabelSeparator*{emdash}{~--- }
\captionsetup[figure]{labelsep=emdash,font=onehalfspacing,position=bottom}
%--------------------------------------

\usepackage{tempora}
\usepackage{amsmath}
\usepackage{color}
\usepackage{listings}
\lstset{
  belowcaptionskip=1\baselineskip,
  breaklines=true,
  frame=L,
  xleftmargin=\parindent,
  language=Python,
  showstringspaces=false,
  basicstyle=\footnotesize\ttfamily,
  keywordstyle=\bfseries\color{blue},
  commentstyle=\itshape\color{purple},
  identifierstyle=\color{black},
  stringstyle=\color{red},
  extendedchars=\true,
}

%--------------------------------------
%			НАЧАЛО ДОКУМЕНТА
%--------------------------------------

\begin{document}

%--------------------------------------
%			ТИТУЛЬНЫЙ ЛИСТ
%--------------------------------------
\begin{titlepage}
\thispagestyle{empty}
\newpage


%Шапка титульного листа
%--------------------------------------
\vspace*{-60pt}
\hspace{-65pt}
\begin{minipage}{0.3\textwidth}
\hspace*{-20pt}\centering

\includegraphics[width=\textwidth]{emblem}
\end{minipage}
\begin{minipage}{0.67\textwidth}\small \textbf{
\vspace*{-0.7ex}
\hspace*{-6pt}\centerline{Министерство науки и высшего образования Российской Федерации}
\vspace*{-0.7ex}
\centerline{Федеральное государственное бюджетное образовательное учреждение }
\vspace*{-0.7ex}
\centerline{высшего образования}
\vspace*{-0.7ex}
\centerline{<<Московский государственный технический университет}
\vspace*{-0.7ex}
\centerline{имени Н.Э. Баумана}
\vspace*{-0.7ex}
\centerline{(национальный исследовательский университет)>>}
\vspace*{-0.7ex}
\centerline{(МГТУ им. Н.Э. Баумана)}}
\end{minipage}
%--------------------------------------

%Полосы
%--------------------------------------
\vspace{-25pt}
\hspace{-35pt}\rule{\textwidth}{2.3pt}

\vspace*{-20.3pt}
\hspace{-35pt}\rule{\textwidth}{0.4pt}
%--------------------------------------

\vspace{1.5ex}
\hspace{-35pt} \noindent \small ФАКУЛЬТЕТ\hspace{80pt} <<Информатика и системы управления>>

\vspace*{-16pt}
\hspace{47pt}\rule{0.83\textwidth}{0.4pt}

\vspace{0.5ex}
\hspace{-35pt} \noindent \small КАФЕДРА\hspace{50pt} <<Теоретическая информатика и компьютерные технологии>>

\vspace*{-16pt}
\hspace{30pt}\rule{0.866\textwidth}{0.4pt}
  
\vspace{11em}

\begin{center}
\Large {\bf Лабораторная работа № 10} \\ 
\large {\bf<<Реализация метода А.А. Карацубы>>} \\ 
\large {\bf по курсу <<Численные методы линейной алгебры>>} \\ 
\end{center}\normalsize

\vspace{8em}


\begin{flushright}
  {Студент группы ИУ9-71Б Лысенко В. В.\hspace*{15pt} \\
  \vspace{2ex}
  Преподаватель Посевин Д. П.\hspace*{15pt}}
\end{flushright}

\bigskip

\vfill
 

\begin{center}
\textsl{Москва 2024}
\end{center}
\end{titlepage}
%--------------------------------------
%		КОНЕЦ ТИТУЛЬНОГО ЛИСТА
%--------------------------------------

\renewcommand{\ttdefault}{pcr}

\setlength{\tabcolsep}{3pt}
\newpage
\setcounter{page}{2}

\section{Цель работы}\label{Sect::task}
\par
\renewcommand{\labelenumii}{\arabic{enumi}.\arabic{enumii}.}
Реализовать метод А.А. Карацубы для перемножения длинных чисел.

\section{Задачи}\label{Sect::task}
\par
\renewcommand{\labelenumii}{\arabic{enumi}.\arabic{enumii}.}
\begin{enumerate} 
    \item Сгенерировать пары чисел размерности $2^{n}$, где n меняется от 1 до 15 (условно).
    \item Перемножить числа обычным методом и методом Карацубы.
    \item Сравнить время работы алгоритмов.
\end{enumerate}

\pagebreak
\section{Практическая реализация}

\par
Код программы:
\begin{lstlisting}
using Random
using Plots

function random_number_array(n)
    return rand(0:9, n)
end


function standard_multiplication(x, y)
    n = length(x)
    result = zeros(Int, 2n)
    for i in 1:n
        carry = 0
        for j in 1:n
            result[i + j - 1] += x[i] * y[j] + carry
            carry = div(result[i + j - 1], 10)
            result[i + j - 1] = mod(result[i + j - 1], 10)
        end
        result[i + n] += carry
    end
    return result
end


function karatsuba(x, y)
    n = length(x)

    if n < 4
        return standard_multiplication(x, y)
    end

    m = div(n, 2)

    x1, x0 = x[1:m], x[m+1:end]
    y1, y0 = y[1:m], y[m+1:end]

    z2 = karatsuba(x1, y1)
    z0 = karatsuba(x0, y0)

    x1x0 = [x1[i] + x0[i] for i in 1:m]
    y1y0 = [y1[i] + y0[i] for i in 1:m]

    z1 = karatsuba(x1x0, y1y0)

    for i in 1:length(z2)
        z1[i] -= z2[i] + z0[i]
    end

    result = zeros(Int, 2n)

    for i in 1:length(z2)
        result[i] += z2[i]
        result[i + m] += z1[i]
        result[i + 2m] += z0[i]
    end

    return result
end

sizes = [2^i for i in 1:15]
std_times = Float64[]
karatsuba_times = Float64[]

for n in sizes
    x_arr = random_number_array(n)
    y_arr = random_number_array(n)

    std_time = @elapsed standard_multiplication(x_arr, y_arr)
    push!(std_times, std_time)

    karatsuba_time = @elapsed karatsuba(x_arr, y_arr)
    push!(karatsuba_times, karatsuba_time)
end

plot(sizes, std_times, label="Обычное умножение", xlabel="Количество разрядов (2^n)", ylabel="Время, с", legend=:topright, title="Сравнение сложности умножения") 
plot!(sizes, karatsuba_times, label="Алгоритм Карацубы", linestyle=:dash)
\end{lstlisting}

\section{Результаты}
\par
Результат выполнения программы:
\begin{center}
\includegraphics[scale=0.7]{g1}
\end{center}
\par


\section{Выводы}
В ходе выполнения лаборатороной работы был реализован метод А.А. Карацубы для перемножения чисел большой размерности. Сравнение с обычным умножением показало, что метод Карацубы действительно отрабатывает гораздо быстрее на больших показателях длины чисел.
\end{document}
